% ===================================================================
%  કોષ્ટક નં. ૧૪ — રસ્તો. --- વેપારીઓ.
%
%  Traduction, faite en 2026, en gujarati d'aujourd'hui, sur l'ido de
%  texto/io. Regles de la colonne : voir 10-kostak-01.tex. Une seule
%  numerotation. 42 blocs, 97 renvois — le tableau des metiers.
%
%  LA VITRINE DU BIJOUTIER (37) EST LA SEULE BOUTIQUE DE CETTE RUE
%  DONT LE GUJARAT TIENNE AUJOURD'HUI LE COMMERCE. Le livret dessine
%  un horloger-bijoutier de sous-prefecture francaise en 1926 et
%  enumere ce qu'il y a dans sa vitrine : « હીરા » les diamants,
%  « મોતી » les perles, « ઝવેરાત » les bijoux. Surat taille et polit
%  aujourd'hui environ neuf diamants sur dix dans le monde, et Cambay
%  et Surat ont enfile les perles de l'ocean Indien pendant des
%  siecles. Les mots ne sont donc pas ceux d'un vocabulaire de luxe
%  appris dans un livre : ce sont les mots d'une industrie qui fait
%  vivre un million de personnes a une heure de route de la ou cette
%  langue se parle. Le bijoutier est « ઝવેરી » et l'orfevre
%  « સોની » — deux noms qui sont aussi des noms de famille, et Zaveri
%  Bazaar est une adresse.
%
%  MAIS « વેપારી » A ETE PREFERE A « વાણિયો », ET IL FAUT DIRE
%  POURQUOI. Le mot que tout le monde connait pour marchand est
%  વાણિયો, le Bania — du sanskrit वणिज् —, celui-la meme qui est passe
%  par le portugais dans l'anglais banyan, et qui a fini par nommer
%  l'arbre sous lequel ces marchands-la tenaient boutique. C'aurait
%  ete le quatrieme mot du releve a remonter le courant apres કુલી au
%  tableau 12 et શતરંજ au tableau 13. Mais વાણિયો est un nom de caste
%  avant d'etre un nom de metier, et le commercant (65) de ce
%  tableau-ci est un drapier francais : on ecrit « વેપારી », qui ne
%  dit que le metier. LE TABLEAU 7 AVAIT NOTE QUE LES METIERS SONT DES
%  CASTES ; C'EST ICI QUE CELA COUTE QUELQUE CHOSE.
%
%  LA FONTAINE DE PIERRE (78) EST « પરબ », ET LE MOT PORTE UN DON. Un
%  પરબ n'est pas un equipement municipal : c'est un point d'eau
%  qu'un particulier fait installer et entretient a ses frais pour que
%  les passants boivent — une fondation pieuse, au meme titre qu'un
%  puits ou un abreuvoir. La planche grave une borne-fontaine de rue.
%  Le mot est exact pour l'objet et ajoute une intention que le
%  francais n'a pas.
%
%  ET L'HARMONIUM (alinea 4) EST LE CONTRAIRE DE TOUS LES REFUS DE
%  CETTE COLONNE. Dans la vitrine du marchand de musique il y a des
%  orgues, des pianos et des harmoniums ; les deux premiers sont
%  transcrits. Le troisieme est « હાર્મોનિયમ » et c'est un mot
%  gujarati plein : l'instrument est arrive avec les missionnaires au
%  XIXe siecle, All India Radio l'a banni de 1940 a 1971 parce qu'il
%  etait etranger, et il est aujourd'hui l'instrument de tout bhajan
%  et de tout ghazal. SEPT FOIS CETTE COLONNE A REFUSE UNE PLANTE OU
%  UN PLAT PARCE QUE LE MOT EXISTAIT SANS QUE LA CHOSE SOIT LA MEME ;
%  VOICI LE CAS INVERSE — la chose est venue, elle est restee, et le
%  mot est a nous.
%
%  Enfin le cercueil (alinea 2). « શબપેટી » existe, et il sert : les
%  Parsis, les chretiens et les musulmans du Gujarat enterrent. Mais
%  la majorite incinere sur une « નનામી », une civiere ouverte, et il
%  n'y a pas de biere. La planche montre un corbillard et un cercueil.
%  On l'ecrit, sans note — comme le cochon du tableau 7 et le menu du
%  tableau 13.
% ===================================================================

%%K t14-apar-1 apar
\VUsaut{4.0mm}
\VUtitre{15.0pt}{60}{કોષ્ટક નં. ૧૪}
\VUsaut{3.0mm}

%%K t14-tit-1 sub
\VUpk{11.4pt}{40}{રસ્તો. --- વેપારીઓ.}
\VUsaut{3.0mm}

%%K t14-01-1 p
1. --- મારો મિત્ર યોહાનેસ, જે ગામડામાં રહે છે, મારા કુટુંબમાં ત્રણ
દિવસ ગાળવા આવ્યો. \VUgras{ત્યાં સુધી} તેને મોટું શહેર જોવાની તક મળી ન
હતી ; એટલે અમે અમારા બધા દિવસ લટાર મારવામાં ગાળીએ છીએ, અને તે જે
જાણતો નથી તે હું તેને સમજાવું છું.

%%K t14-02-1 p
2. --- પહેલાં અમે \VUgras{વેધશાળા} \textsuperscript{(1)} જોવા ગયા, જેમાં
તેને ખૂબ રસ પડ્યો. પાછા ફરતાં એ \VUgras{રસ્તા} \textsuperscript{(3)}
પરથી, જ્યાં \VUgras{તુર્કી હમામ} \textsuperscript{(2)} છે, અમને એક
\VUgras{સ્મશાનયાત્રા} \textsuperscript{(4)} મળી. \VUgras{અંતિમયાત્રાની}
આગળ \VUgras{પાદરીઓ} ચાલતા હતા, \VUgras{શબવાહિનીની} આગળ, જેમાં
\VUgras{શબપેટી} મુકાયેલી હતી. ઘણા મિત્રો \VUgras{શોકગ્રસ્ત} કુટુંબની
સાથે હતા.

%%K t14-02-2 p
સરઘસ પસાર થયા પછી અમે મોટા \VUgras{બિયરઘર} \textsuperscript{(5)} આગળ
રસ્તો ઓળંગ્યો, જ્યાં ઘણા \VUgras{ગ્રાહકો} અગાસી પર \VUgras{બિયર}
પીતા હતા, કાચના \VUgras{છજા} \textsuperscript{(6)} નીચે આશરો લઈને.

%%K t14-03-1 p
3. --- \VUgras{દારૂના વેપારીની} \textsuperscript{(7)} દુકાન અને
\VUgras{સંગીતની દુકાનવાળાની} \textsuperscript{(10)} દુકાન વચ્ચે મુકાયેલા
\VUgras{ઢાલચિહ્નથી} યોહાનેસ ખૂબ નવાઈ પામ્યો. તેણે મને તેનો અર્થ
પૂછ્યો ; મેં જવાબ આપ્યો કે તે અંગ્રેજ \VUgras{વાણિજ્યદૂત કચેરી}
\textsuperscript{(8)} બતાવે છે, અને તેને ઝરૂખા પરના
\VUgras{વાણિજ્યદૂત} \textsuperscript{(9)} બતાવ્યા.

%%K t14-tit-2 sub
\VUpk{10.2pt}{40}{દુકાનોની તપાસ.}
\VUsaut{3.0mm}

%%K t14-04-1 p
4. --- સ્વાભાવિક રીતે અમે \VUgras{વાજિંત્રોની}, \VUgras{હાર્મોનિયમની},
\VUgras{ઑર્ગનની} અને \VUgras{પિયાનોની} માંડણી આગળ થોભ્યા ; અમે
\VUgras{સ્વરલિપિનાં} અને પિયાનો માટેની \VUgras{સંગીતની ચોપડીઓનાં}
ચિત્રવાળાં પૂંઠાં જોયાં, અને સોનેરી લાકડાની \VUgras{લાયર} જોઈને વાહ
પોકારી, જે નિશાની તરીકે વપરાય છે.

%%K t14-05-1 p
5. --- \VUgras{ઝાડુવાળાઓએ} \textsuperscript{(11)} અને
\VUgras{સફાઈગાડીએ} \textsuperscript{(12)} ઉડાડેલી ધૂળનાં વાદળોએ અમને
\VUgras{રાચરચીલાના વેપારીના} \textsuperscript{(13)} સુંદર
\VUgras{રાચરચીલા} આગળથી અને \VUgras{કલઈગરની} \textsuperscript{(14)}
\VUgras{ચૂલા} અને \VUgras{ફાનસની} માંડણી આગળથી ઝડપથી પસાર થવાની ફરજ
પાડી.

%%K t14-06-1 p
6. --- વળી, આ જગ્યાએ અમારે ફૂટપાથ છોડવી પડી, કેમ કે તે પોટલાંથી
ભરાયેલી હતી, જે એક \VUgras{માલગાડામાંથી} \textsuperscript{(16)}
ઉતારાતાં હતાં, હમણાં જ ભાડે લેવાયેલી એક \VUgras{દુકાનમાં}
\textsuperscript{(15)} મૂકવા માટે.

%%K t14-06-2 p
એથી અમે \VUgras{ગાડી બનાવનારની} \textsuperscript{(17)} સુંદર ગાડીઓ ન
જોઈ શક્યા, પણ એથી અમે \VUgras{પૅકિંગ કરનારને} \textsuperscript{(19)}
જોવા થોભતાં અટક્યા નહીં, જે પોતાના પડોશી — \VUgras{સાઇકલ અને મોટરના
વેપારી} \textsuperscript{(18)} — માટે એક પેટી બનાવતો હતો. પછી, થોડું
મોડું થઈ ગયું હતું એટલે, અમે ઘરે પાછા ફર્યા.

%%K t14-tit-3 sub
\VUpk{10.2pt}{40}{શહેરમાં લટાર.}
\VUsaut{3.0mm}

%%K t14-07-1 p
7. --- બીજા દિવસે મારી માએ યોહાનેસને સૂચવ્યું કે તેનો ફોટો પડાવીએ,
અને પછી તેને \VUgras{ફોટોગ્રાફરને} \textsuperscript{(20)} ત્યાં લઈ
જવાનું કામ મને સોંપ્યું. જમ્યા પછી અમે એ કલાકારના \VUgras{સ્ટુડિયોમાં}
ગયા, જ્યાં અમારે ઠીક ઠીક લાંબી રાહ જોવી પડી. વખત કાઢવા અમે દીવાલે
લટકતી \VUgras{છબીઓ} \textsuperscript{(22)} જોઈ.

%%K t14-07-2 p
અમારો વારો આવ્યો ત્યારે કામ લાંબું ચાલ્યું નહીં, અને મારા મિત્રે
\VUgras{કૅમેરા} \textsuperscript{(21)} આગળ ઊભા રહેવાનું પૂરું કર્યું કે
તરત અમે નીચે ઊતર્યા.

%%K t14-08-1 p
8. --- પહેલા માળે અમે \VUgras{હજામના} \textsuperscript{(23)} બારણામાં
થઈને જોયું, જે ખુલ્લું છે, કે તેનો છોકરો એક ગ્રાહકના વાળ કાપતો હતો.

%%K t14-08-2 p
રસ્તા પર આવીને અમે \VUgras{તમાકુની દુકાન} \textsuperscript{(24)} આગળથી
પસાર થયા. લાલ \VUgras{ફાનસ} \textsuperscript{(25)} અને \VUgras{મોટી
ચલમ}, જે \VUgras{નિશાની} \textsuperscript{(26)} તરીકે વપરાય છે, તેનાથી
યોહાનેસને એટલી મજા પડી કે તે બે ઘરડા સાહેબો સાથે અથડાયો, જે
\VUgras{છીંકણી સૂંઘતાં} \textsuperscript{(27)} વાત કરતા હતા.

%%K t14-tit-4 sub
\VUpk{10.2pt}{40}{મીઠાઈની દુકાન.}
\VUsaut{3.0mm}

%%K t14-09-1 p
9. --- \VUgras{શરાફના} \textsuperscript{(29)} \VUgras{શોકેસમાં}
માંડેલી બધા દેશની \VUgras{નોટો} અને \VUgras{સિક્કા} થોડી વાર અમારું
ધ્યાન ખેંચી રહ્યાં, પણ નાસ્તાનો સમય થયો હતો એટલે અમે વધારે થોભ્યા વિના
\VUgras{મીઠાઈવાળાને} \textsuperscript{(31)} ત્યાં પેઠા.

%%K t14-10-1 p
10. --- દુકાનમાં જે બધી \VUgras{મીઠાઈઓ} હતી — \VUgras{કેક, મધની કેક,
બિસ્કિટ} અને જાતજાતની \VUgras{પીપરમીટ} — તે જોઈને અમે સહેલાઈથી પસંદ
કરી શક્યા નહીં. છેવટે યોહાનેસે \VUgras{મલાઈની કેક} પસંદ કરી, અને મેં
માત્ર એક નાની \VUgras{ચેરીની ટાર્ટ} લીધી, કેમ કે \VUgras{ગળ્યું મારા
દાંતને દુખાડે છે}, અને મને \VUgras{દાંતના ડૉક્ટર} \textsuperscript{(30)}
પાસે વારે વારે જવાનું જરાય ગમતું નથી.

%%K t14-tit-5 sub
\VUpk{10.2pt}{40}{ટાઇપરાઇટર પર લખવું.}
\VUsaut{3.0mm}

%%K t14-11-1 p
11. --- મારા મિત્રને \VUgras{ટાઇપરાઇટર} બતાવવા, જેની વાત તેણે સાંભળી
હતી પણ જે તેણે જોયું ન હતું, અમે એક \VUgras{કચેરીમાં}
\textsuperscript{(32)} પેઠા, જે મારા \VUgras{પિતરાઈ ભાઈની}
\textsuperscript{(34)} છે. અમે તેમને પોતાના \VUgras{સેક્રેટરીને}
\textsuperscript{(35)} પત્રો લખાવતા જોયા, જે \VUgras{લઘુલિપિકાર} છે. પત્ર પૂરો થાય કે તરત એક
\VUgras{ટાઇપિસ્ટ} \textsuperscript{(33)} તેને પોતાના યંત્ર પર ઉતારી
લેતો. મારા સગાના કામમાં ખલેલ ન પાડવાની ઇચ્છાથી અમે તેમની પાસે માત્ર
થોડી જ ક્ષણ રહ્યા.

%%K t14-12-1 p
12. --- મારા પિતરાઈની કચેરીની નીચે એક મોટા \VUgras{ઘડિયાળી અને ઝવેરી}
\textsuperscript{(37)} રહે છે, જે વળી \VUgras{સોની} પણ છે. તેમની ભવ્ય
દુકાનમાં ઘણાં \VUgras{ઘડિયાળ, લોલકવાળાં ઘડિયાળ, ખિસ્સાઘડિયાળ},
\VUgras{સાંકળીઓ} અને કીમતી \VUgras{ઝવેરાત} છે — દાખલા તરીકે
\VUgras{મોતીના હાર}, સોનાની \VUgras{બંગડીઓ}, \VUgras{કાનની બુટ્ટીઓ}
અને \VUgras{વીંટીઓ}, જે \VUgras{રત્નો} અને \VUgras{હીરાથી} જડેલી છે.
એક શોકેસ ખાસ \VUgras{સોનાની વસ્તુઓ} માટે રાખેલો છે, અને તેમાં ચાંદીનાં
\VUgras{ખાવાનાં વાસણોનો} મોટો સંગ્રહ છે.

%%K t14-13-1 p
13. --- રસ્તાના ખૂણે વળતાં મેં જોયું કે ઘરના \VUgras{નંબર} પાસે
મુકાયેલી \VUgras{પાટી} \textsuperscript{(36)} બદલાઈ ગઈ છે. હું મોટેથી
મારી નોંધ કહેવા જતો હતો ત્યાં \VUgras{લશ્કરી બૅન્ડના} આનંદી સૂર
સંભળાયા. અમે પલટણ સામે દોડ્યા, એ વિચાર્યા વિના કે \VUgras{તે}
\VUgras{ટોપીવાળાની} \textsuperscript{(39)} અને \VUgras{મોજાંવાળાની}
\textsuperscript{(40)} દુકાનો આગળથી પસાર થવાની છે, અને એ કે તેની કૂચ
જોવા માટે અમે \VUgras{ચશ્માંવાળાના} \textsuperscript{(38)} શોકેસ આગળ
ઊભા રહીને પણ એટલી જ સારી જગ્યાએ હોત, જે \VUgras{નાક પર બેસતાં
ચશ્માં, કાનવાળાં ચશ્માં} અને \VUgras{હાથે પકડવાનાં ચશ્માંથી} ભરેલો
છે.

%%K t14-tit-6 sub
\VUpk{10.2pt}{40}{કૂચનું સરઘસ.}
\VUsaut{3.0mm}

%%K t14-14-1 p
14. --- \VUgras{પલટણની} \textsuperscript{(41)} આગળ \VUgras{ઢોલનો
આગેવાન} \textsuperscript{(42)} ચાલતો હતો, તેની પાછળ \VUgras{ઢોલીઓ}
\textsuperscript{(43)}, \VUgras{બ્યુગલવાળા} \textsuperscript{(44)} અને
આખું \VUgras{બૅન્ડ}. \VUgras{જનરલ} \textsuperscript{(45)}, જે પોતાના
મદદનીશ સાથે ચોકમાં હતા, તે \VUgras{ફુવારા} \textsuperscript{(49)}
પાસે થોભ્યા હતા, જે \VUgras{પાણીના હોજનો} \textsuperscript{(48)} છે,
\VUgras{કૂચ} જોવા માટે.

%%K t14-14-2 p
\VUgras{વડા મથકના અફસરોએ} \textsuperscript{(46)}, \VUgras{કર્નલે} અને
\VUgras{લેફ્ટનન્ટ કર્નલે} તલવારથી તેમને સલામ કરી. \VUgras{મેજરો}
પોતાની \VUgras{બટાલિયનની} \textsuperscript{(47)} આગળ હતા અને દરેક
\VUgras{કપ્તાન} પોતાની \VUgras{ટુકડીને} હુકમ કરતો હતો, જ્યારે
\VUgras{લેફ્ટનન્ટો, સબ-લેફ્ટનન્ટો} અને \VUgras{નાયબ અફસરો} પોતાના
માણસો પાસે પોતાની રૂઢ જગ્યાએ હતા.

%%K t14-15-1 p
15. --- ઘણા વટેમાર્ગુ \VUgras{ચાર રસ્તા} \textsuperscript{(50)} પર
ટોળે વળ્યા હતા ; વેપારીઓ, \VUgras{છત્રીવાળો} \textsuperscript{(51)},
\VUgras{રંગારો} \textsuperscript{(53)}, \VUgras{હથિયારવાળો}
\textsuperscript{(54)} \VUgras{સિપાહીઓને} જોવા બહાર નીકળ્યા. બધાં
વાહન — \VUgras{ભાડાની ઘોડાગાડીઓ} \textsuperscript{(52)},
\VUgras{ખાનગી ઘોડાગાડીઓ} \textsuperscript{(57)} અને \VUgras{પાણી
છાંટવાનું પીપ} \textsuperscript{(56)} પણ — ફૂટપાથની બાજુમાં ઊભાં રહી
ગયાં.

%%K t14-tit-7 sub
\VUpk{10.2pt}{40}{રસ્તા પરનાં દૃશ્યો.}
\VUsaut{3.0mm}

%%K t14-15-2 p
એક \VUgras{વેપારનો ફરતો પ્રતિનિધિ} \textsuperscript{(55)}, હાથમાં
પોતાની \VUgras{નમૂનાની પેટીઓ} લઈને, ઝડપથી રસ્તો ઓળંગી ગયો, અને
\VUgras{ઉપલી બેઠક} \textsuperscript{(59)} પર બેઠેલા લોકો ---
\VUgras{ઓમ્નિબસની} \textsuperscript{(58)} બેઠક --- જોવાની મજા લેવા
પાછળ ફર્યા.
બિચારો \VUgras{ફરતો ગવૈયો} \textsuperscript{(60)}, પોતાના \VUgras{હાથે
ફેરવવાના વાજા} \textsuperscript{(61)} સાથે, પોતાનું \VUgras{ગીત}
અધવચ્ચે બંધ કરવું પડ્યું, કેમ કે ઢોલના ધડાકાએ તેનો અવાજ ઢાંકી દીધો.

%%K t14-16-1 p
16. --- પલટણ \VUgras{કાપડની દુકાન} \textsuperscript{(64)} આગળ પહોંચી
ત્યારે \VUgras{સીવણ કરનારી બહેનો} \textsuperscript{(63)} અને
\VUgras{દરજીઓ} \textsuperscript{(62)} પોતાના \VUgras{સંચા} અને પોતાનું
કામ છોડીને બારીમાંથી જોવા લાગ્યાં. \VUgras{વેપારીએ}
\textsuperscript{(65)}, જેણે હમણાં જ પોતાની \VUgras{માંડણીએ} નવા
\VUgras{પોશાક} \textsuperscript{(66)} ગોઠવ્યા હતા, તેણે \VUgras{કાપડના તાકા}
\textsuperscript{(67)} અને \VUgras{કપડાં} અંદર મુકાવવાં પડ્યાં, જે
ફૂટપાથ પર \VUgras{ગટરના ઢાંકણા} \textsuperscript{(68)} પાસે માંડેલાં
હતાં, રખેને ટોળું તેમને ઢાળી દે ; ઘરડા \VUgras{ફેરિયાને}
\textsuperscript{(69)} સુધ્ધાં, જેની પાસેથી એક રખેવાળ બહેન મોજાંનો ભાવ
કરતી હતી, પોતાનું વેચાણ પૂરું કરવા એક પરસાળમાં પેસવું પડ્યું.

%%K t14-16-2 p
અમે પલટણની સાથે છાવણી સુધી ગયા, \VUgras{અને,} અમારા દિવસથી ખૂબ સંતોષ
પામીને, રાતના જમણના સમયે અમે પાછા ફર્યા.

%%K t14-tit-8 sub
\VUpk{10.2pt}{40}{જુદા જુદા વેપાર અને ધંધા.}
\VUsaut{3.0mm}

%%K t14-17-1 p
17. --- બીજા દિવસે મારા પિતાએ અમને ત્યાંના છાપાનું છાપખાનું જોવા
જવાનું સૂચવ્યું. અમે ઉત્સાહથી હા પાડી.

%%K t14-17-2 p
\VUgras{છાપનારો} વળી \VUgras{પુસ્તકવિક્રેતા} \textsuperscript{(70)}
\VUgras{(= ચોપડી વેચનારો),} \VUgras{પ્રકાશક}, \VUgras{કાગળનો વેપારી}
અને \VUgras{પુસ્તક બાંધનારો} પણ છે. તેણે પહેલા માળે છાપાનું
\VUgras{તંત્રીખાનું} \textsuperscript{(71)} બેસાડ્યું છે, અને
\VUgras{કોતરકામનો ખંડ} પણ, જ્યાં \VUgras{શિલાછાપ કરનારા}
\textsuperscript{(72)} અને \VUgras{તાંબા પર કોતરનારા}
\textsuperscript{(73)} કામ કરે છે, અને છેવટે \VUgras{બીબાં ગોઠવવાનો
ખંડ}, જ્યાં \VUgras{છાપકામના મજૂરો} \textsuperscript{(74)} છે. ખરું
\VUgras{છાપખાનું} \textsuperscript{(75)}, \VUgras{છાપવાનાં યંત્ર} અને
બીજાં જુદાં જુદાં યંત્ર ભોંયતળિયે છે.

%%K t14-tit-9 sub
\VUpk{10.2pt}{40}{યોહાનેસ ઘણા સવાલ પૂછે છે.}
\VUsaut{3.0mm}

%%K t14-18-1 p
18. --- છાપખાનામાંથી નીકળતાં યોહાનેસ ખૂબ નવાઈ પામીને એક નાની કાચની
પેટી આગળ થોભ્યો, જે \VUgras{સૂતર-બટનની દુકાન} \textsuperscript{(77)}
સામે એક થાંભલા પર મુકાયેલી છે, અને પૂછ્યું કે એ શા માટે વપરાય છે.
મારા પિતાએ તેને સમજાવ્યું કે એ \VUgras{આગની ચેતવણીની પેટી}
\textsuperscript{(76)} છે. બાકી શહેરમાં બધું જ મારા મિત્રને માટે
નવાઈભર્યું હતું — સાદા \VUgras{પથ્થરના પરબથી} \textsuperscript{(78)}
અને હલકી \VUgras{માલ ઊંચકવાની ત્રિચક્રીથી}
\textsuperscript{(80)}, જે \VUgras{વરદીવાળો છોકરો}
\textsuperscript{(79)} ચલાવે છે, માંડીને
\VUgras{ટપાલગાડી} \textsuperscript{(81)} સુધી.

%%K t14-19-1 p
19. --- મારા પિતા થોડી વાર અમને છોડીને \VUgras{વેરા ઉઘરાવનાર
અમલદાર} \textsuperscript{(83)} સાથે વાત કરવા ગયા, જે એક \VUgras{વકીલ}
\textsuperscript{(82)} અને એક \VUgras{નોટરી} \textsuperscript{(84)}
સાથે વાત કરવા થોભ્યા હતા ; એ દરમિયાન અમે રસ્તા પરની અવરજવર નજરે
પીછો કરીને મજા કરી. જાતજાતના માણસો અમારી સામેથી પસાર થયા : એક
\VUgras{મીઠાઈવાળાનો છોકરો} \textsuperscript{(85)}, ટોપલામાં ભવ્ય
\VUgras{પેસ્ટ્રી} લઈને, એક \VUgras{તૈયાર કપડાંના વેપારીની}
\textsuperscript{(86)} પાછળ આવ્યો ; એક \VUgras{ફાનસ સળગાવનારો}
\textsuperscript{(87)} એક \VUgras{ગૅસવાળા} \textsuperscript{(88)} સાથે
વાત કરતો હતો ; એક \VUgras{સાધ્વી} \textsuperscript{(89)}, હાથે એક અનાથ
છોકરીને પકડીને, પોતાના મઠમાં પાછી જતી હતી.

%%K t14-tit-10 sub
\VUpk{10.2pt}{40}{ઘરે પાછા ફરતાં.}
\VUsaut{3.0mm}

%%K t14-20-1 p
20. --- મારા પિતા અમને પાછા આવી મળ્યા એ ઘડીએ યોહાનેસ ખડખડાટ હસી
પડ્યો, બે \VUgras{ભૂંગળું સાફ કરનારાનાં} \textsuperscript{(91)} મેલાં
મોઢાં અને વિચિત્ર ફાટેલાં કપડાં જોઈને, જે તદ્દન મેશથી કાળા હતા.

%%K t14-21-1 p
21. --- મારે મારી મા માટે \VUgras{ફૂલવાળી} \textsuperscript{(92)}
પાસેથી એક છોડ ખરીદવો હતો, પણ મારા પિતાએ મને સમજાવ્યું કે એ ઊંચકીને લઈ
જવા બહુ ભારે પડશે અને \VUgras{સંતરાં} લેવાં વધારે સારું રહેશે. અમે
\VUgras{વેચનારી બહેન} \textsuperscript{(94)} પાસે ગયા, અને જ્યારે
\VUgras{સંભાળ રાખનારી બહેને} \textsuperscript{(93)}, જે પોતાના છોકરા
માટે પસંદ કરતી હતી, પોતાની ખરીદી પૂરી કરી, ત્યારે અમે અમારો વારો
લીધો.

%%K t14-22-1 p
22. --- ઘરે પાછા ફરતાં અમને કેટલાક ઓળખીતા મળ્યા : મારા કુટુંબની એક
ઘરડી બહેનપણી, જે પોતાની \VUgras{સાથે રહેનારી બહેનના}
\textsuperscript{(95)} હાથ પર ટેકો દેતી હતી ; પુલ અને રસ્તાના
\VUgras{ઇજનેર} \textsuperscript{(96)} ; અને મારી એક પિતરાઈ બહેન પોતાની
\VUgras{આયા} \textsuperscript{(98)} સાથે.

%%K t14-22-2 p
પછી અમને મારી માની \VUgras{ટોપી બનાવનારી બહેન} \textsuperscript{(97)}
મળી અને શહેરમાં નવા એવા બે \VUgras{સાધુઓ} \textsuperscript{(99)}, જે
મારા પિતાને સ્ટેશનનો રસ્તો પૂછવા અમારી પાસે આવ્યા.

\VUsaut{6.0mm}
\VUfilet{22mm}
